\documentclass[10pt,a4paper]{amsart}
\usepackage[margin=19mm]{geometry}
\usepackage{amsmath,amssymb,mathtools}
\usepackage[scr=rsfs,cal=euler]{mathalfa}
\usepackage{xfrac}
\usepackage[unicode,hidelinks,bookmarks=false]{hyperref}
\newcommand{\ie}{i.e.\ }
\newcommand{\cf}{cf.\ }
\newcommand{\resp}{respectively\ }
\newcommand{\Subsection}{\subsection}
\providecommand{\nobreakdash}{\nobreak\hbox{-}}
\setlength{\emergencystretch}{2em}
\multlinegap0pt
\usepackage{bm}
\usepackage{bbm}
\usepackage{dsfont}
\usepackage{xspace}
\usepackage{cancel}
\usepackage{mathtools}
\usepackage{amsthm}
\makeatletter
\@ifundefined{theorem}{\newtheorem{theorem}{Theorem}}{}
\@ifundefined{lemma}{\newtheorem{lemma}[theorem]{Lemma}}{}
\@ifundefined{proposition}{\newtheorem{proposition}[theorem]{Proposition}}{}
\makeatother
\makeatletter
\expandafter\ifx\csname acknowledgement\endcsname\relax
\newenvironment{acknowledgement}[1][Acknowledgement
]{\begin{trivlist} \item[\hskip \labelsep {\bfseries
#1}]}{\end{trivlist}}
\fi
\expandafter\ifx\csname example\endcsname\relax
\newenvironment{example}[1][Example]{\begin{trivlist}
\item[\hskip \labelsep {\bfseries #1}]}{\end{trivlist}}
\fi
\makeatother
\title[Stability of Inverse Resonance Problem on the Half Line]{Stability of Inverse Resonance Problem on the Half Line}
\author{Lung-Hui Chen}
\date{Source: arXiv:2508.11103v1 (CC BY 4.0)}
\begin{document}
\maketitle

\noindent\emph{Source and licence.} Stability of Inverse Resonance Problem on the Half Line. Version arXiv:2508.11103v1 (CC BY 4.0), distributed under the Creative Commons Attribution 4.0 International licence, as recorded in the arXiv licence metadata for this exact version and shown on the abstract page https://arxiv.org/abs/2508.11103v1. Full rights notice: licenses/arxiv-2508.11103-CC-BY-4.0.txt.\par\smallskip

\noindent\textit{Editorial scope.} Counted unit: the Proposition whose source label is \texttt{D} in the article (source0.tex lines 311--315, source label \texttt{D}), delivered together with the article's own complete original proof of it (source0.tex lines 316--337, one \texttt{proof} environment of 22 lines). No mathematics is written, repaired or re-derived: statement and proof are the article's own text line for line. Required context retained uncounted: VV (lines 193--199); sum (lines 209--211); mu (lines 247--249); Dickson (lines 292--307); 2.9 (lines 302--304). Every definition, display and label of those lines is retained unchanged. Omitted from this package and not redistributed: 1--192, 200--208, 212--246, 250--291, 305--310, 338--508. Nothing inside a retained excerpt is omitted or altered: every retained body is delivered exactly as the article's lines are written, so no omission receipt of any kind accompanies this package. Citations: the retained excerpts carry no citation command at all, so no bibliography entry of the article is transcribed and no reference is redistributed. Reference adaptation: none was needed and none was made; every reference inside the delivered unit resolves inside it, so no unresolved reference or citation is produced. Printed slips kept literal: no printed word, symbol, hypothesis or number of the counted statement or of its proof was altered. Layout and class: the article's own class is \texttt{article} with packages including fullpage, amsfonts, amsmath, amssymb, amsthm, graphicx, sectsty, authblk; the wrapper is the lane's pinned \texttt{amsart} editorial wrapper at 10pt on A4, which restates the theorem-like environments the retained text needs and adds the article's own macro definitions that the retained text needs (none), with extra wrapper packages (bm, bbm, dsfont, xspace, cancel, mathtools). The delivered numbering is the unit's own. Assets: no retained span contains a figure environment, a graphics-inclusion command, a PDF-page inclusion, a table or a TikZ picture; no image file is copied into this package, the package declares no asset dependency and no image file is redistributed with it. Licence: version arXiv:2508.11103v1 of this article is distributed under the Creative Commons Attribution 4.0 International licence, as recorded in the arXiv licence metadata for this exact version and shown on the abstract page https://arxiv.org/abs/2508.11103v1. A full rights notice is delivered at licenses/arxiv-2508.11103-CC-BY-4.0.txt. Characters: every retained excerpt is pure ASCII, so the delivered engine, XeLaTeX with its default Latin Modern text font, reports no missing character.
\begin{lemma}
The following asymptotic holds for real $k$:
\begin{eqnarray}\label{VV}
4z^{2}\hat{V}(2z)\hat{V}(-2z)=1+o(\frac{1}{z}),
\end{eqnarray}
in which all little o-terms are analytic.
\end{lemma}

\begin{equation}\label{sum}
f(z)=\sum_{j=1}^nA_jz^{m_j}[1+\epsilon(z)]e^{\omega_jz},\,z\in\mathbb{C},
\end{equation}

\begin{equation} \label{mu}
\mu_{kj}:=\frac{m_{kj}-m_{kj+1}}{(\omega_{kj}-\omega_{kj+1})e_k},
\end{equation}

\begin{theorem}[Dickson]\label{Dickson}
Let $f(z)$ be given as in~(\ref{sum}). Then, there exists $H>0$
such that \begin{enumerate}
               \item all but a finite number of zeros of $f$ of modulus
greater than $r_0$ are in $\cup_{k,j}V_{kj}$;
               \item
for each pair
of positive reals $\epsilon$ and $s_0$, there exists an
$\alpha_0=\alpha_0(\epsilon,s_0)$ such that whenever
$\alpha\geq\alpha_0$ and $s\geq s_0$,
\begin{equation}\label{2.9}
\big|N_f(R_{kj}(\alpha,s,H))-s|\omega_{kj+1}-\omega_{kj}|/(2\pi)\big|<n_{kj}-1+\epsilon,
\end{equation}
in which $n_{kj}$ is given in~(\ref{mu}).
\end{enumerate}
\end{theorem}

\begin{proposition}\label{D}
Let $\{z_{n}^{j}\}_{n=1}^{\infty}$ be the point sequence of zeros of $\hat{V}^{j'}(2z)\hat{V}^{j'}(-2z)$ such that $|z_{1}^{j'}|<|z_{2}^{j'}|<|z_{3}^{j'}|<\cdots$, $j'=1,2.$ For $0\leq\delta\ll 1$ and $0\leq\sup_{n}|z^{1}_{n}-z^{2}_{n}|<\delta$, the difference of numbers of $\{z^{1}_{n}\}$ and $\{z^{2}_{n}\}$ inside  \begin{eqnarray*}
S_{R}:=\{z\in\mathbb{C}|\,0\leq\Re z\leq R,\,-K\leq\Im z\leq  K\},
\end{eqnarray*} is uniformly bounded.
\end{proposition}

\begin{proof}
We apply Theorem \ref{Dickson}. For our exponential polynomial 
$16z^{4}\hat{V}^{j'}(2z/i)\hat{V}^{j'}(-2z/i)$ in~(\ref{VV}), we deduce
\begin{eqnarray*}
&&k=1,\\
&&e_{1}=0,\\
&&\omega_{11}=-2,\,\omega_{12}=0,\,\omega_{13}=2;\\
&&m_{12}=m_{23}=0;\\
&&\mu_{12}=\mu_{23}=0;\\
&&n_{11}=n_{12}=2.
\end{eqnarray*}
Therefore, we use~(\ref{2.9}) to obtain
\begin{equation}\nonumber
\big|N_{\hat{V}^{^{j'}}(2z)\hat{V}^{^{j'}}(-2z)}(R_{1j}(\alpha,s,K))-s|\omega_{1j+1}-\omega_{1j}|/(2\pi)\big|<1+\epsilon.
\end{equation} 
Let us choose $\alpha\gg1$ large and $s=2\pi/|\omega_{1j+1}-\omega_{1j}|$ so that
\begin{equation}\label{2.11}
\big|N_{\hat{V}^{^{j'}}(2z)\hat{V}^{^{j'}}(-2z)}(R_{1j}(\alpha,2\pi/|\omega_{1j+1}-\omega_{1j}|,K))-1\big|<1+\epsilon,\,j'=1,2,\,j=1,2.
\end{equation} 
Hence, for small $0\leq\delta\ll1$, we have $z^{1}_{n}$ and $z^{2}_{n}$ falling in the same rectangle for large $n$.
This proves the proposition.
\end{proof}

\end{document}
